% ============================================================
%  Admissibility Before Truth
%  On the Domain of Truth Predicates
% ============================================================
%  Engine: LuaLaTeX
%  Font:   TeX Gyre Pagella
%  Third paper in the repair-theoretic program.
%  Companion papers:
%    [Flyxion2026a] Observability Is Restorability
%    [Flyxion2026b] Erasure Is Deformation
% ============================================================

\documentclass[12pt, oneside]{article}

% ── Fonts & encoding ────────────────────────────────────────
\usepackage{fontspec}
\setmainfont{TeX Gyre Pagella}
\usepackage{unicode-math}
\setmathfont{TeX Gyre Pagella Math}

% ── Layout ──────────────────────────────────────────────────
\usepackage[
  top=1.25in, bottom=1.25in,
  left=1.35in, right=1.35in
]{geometry}
\usepackage{setspace}
\setstretch{1.15}
\setlength{\parskip}{0.55em}
\setlength{\parindent}{1.4em}

% ── Mathematics ─────────────────────────────────────────────
\usepackage{amsmath, amsthm}
\usepackage{mathtools}

% Theorem environments
\theoremstyle{plain}
\newtheorem{theorem}{Theorem}[section]
\newtheorem{proposition}[theorem]{Proposition}
\newtheorem{lemma}[theorem]{Lemma}
\newtheorem{corollary}[theorem]{Corollary}

\theoremstyle{definition}
\newtheorem{definition}[theorem]{Definition}
\newtheorem{example}[theorem]{Example}
\newtheorem{axiom}[theorem]{Axiom}

\theoremstyle{remark}
\newtheorem{remark}[theorem]{Remark}

% ── Notation shorthands ─────────────────────────────────────
\DeclareMathOperator{\Rest}{Rest}
\DeclareMathOperator{\Obs}{Obs}
\newcommand{\CR}{C_{\!R}}
\newcommand{\AO}{\mathcal{A}}
\newcommand{\Rzero}{\mathcal{R}}
\newcommand{\Dsub}{\mathcal{D}}
\newcommand{\Xsub}{X_S}
\newcommand{\Tsub}{\mathcal{T}}
\newcommand{\Ttilde}{\widetilde{\mathcal{T}}}
\newcommand{\iotamap}{\iota}        % inclusion map Rest(R) -> D
\newcommand{\RestR}{\Rest(\mathcal{R})}
\newcommand{\bA}{\partial\mathcal{A}}
\newcommand{\Undef}{\mathord{\perp}} % unevaluable truth value

% ── Section & title styling ─────────────────────────────────
\usepackage{titlesec}
\titleformat{\section}{\normalfont\large\bfseries}{{\thesection}.}{0.6em}{}
\titleformat{\subsection}{\normalfont\normalsize\bfseries}%
  {{\thesubsection}.}{0.5em}{}

% ── Hyperref (load last) ────────────────────────────────────
\usepackage[
  colorlinks=true,
  linkcolor=black,
  citecolor=black,
  urlcolor=black
]{hyperref}

% ============================================================
\begin{document}

% ── Title block ─────────────────────────────────────────────
\begin{center}
  {\LARGE\bfseries Admissibility Before Truth}\\[0.9em]
  {\large On the Domain of Truth Predicates}\\[1.5em]
  {\normalsize Flyxion}\\[0.3em]
  {\small\itshape Independent Researcher}\\[1.6em]
  {\small\today}
\end{center}

\vspace{0.5em}

% ── Abstract ────────────────────────────────────────────────
\begin{abstract}
\noindent
Standard epistemology treats truth as the primary predicate and asks
which beliefs satisfy it.
This paper argues that truth is not primary but derived: it is a
partial predicate whose domain is determined by a prior structure,
the observer's admissibility manifold.
We prove the \emph{Domain Restriction Theorem}: every well-defined
truth predicate $\Tsub$ factors through the inclusion
$\iotamap: \RestR \hookrightarrow \Dsub$
of the restorable distinction space into the full distinction space,
so that $\operatorname{dom}(\Tsub) \subseteq \RestR$.
Distinctions outside the admissible set are not false; they are
unevaluable, receiving the third value $\Undef$ in a natural
three-valued extension.
The theorem is not a philosophical claim but a structural consequence
of the repair-theoretic framework developed in the companion papers
\cite{Flyxion2026a,Flyxion2026b}: truth-evaluation is itself a repair
process, and can only operate where repair capacity exists.
Consequences are derived for Bayesian inference, machine learning,
legal admissibility, language acquisition, and the Kuhnian account
of scientific revolutions.
In each case, what had appeared to be a domain-specific epistemic
constraint is revealed to be an instance of the same domain
restriction.
The paper is the third in a five-paper program; an appendix
summarises the definitions required from the companion papers so
that the central theorem can be read independently.
\end{abstract}

\vspace{1.2em}
\hrule
\vspace{1.6em}

% ============================================================
\section{The Program and This Paper's Place in It}
\label{sec:intro}
% ============================================================

This paper is the third in a sequence of five, each of which
introduces one genuinely new layer of mathematical structure while
relying on the previous layers rather than rebuilding them.
The architecture is as follows.

The first paper, \emph{Observability Is Restorability}~\cite{Flyxion2026a},
defines the primitive objects: repair operators, restorable
distinctions, the admissibility manifold $\AO(\mathcal{R})$, and
the repair pseudometric.
Its central theorem is that the set of observable distinctions
coincides with the set of restorable distinctions.

The second paper, \emph{Erasure Is Deformation}~\cite{Flyxion2026b},
defines the dynamics of loss: repair-repertoire reduction deforms
the admissibility manifold, generating metric singularities at the
irrestorability boundary.
Its central object is the pullback deformation tensor
$E_t = \iota_t^* g_0 - g_t$ measuring how erasure distorts the
geometry of what remains.

The present paper establishes the logical consequences: truth is a
partial predicate whose domain is the admissible distinction space.
The central theorem is the Domain Restriction Theorem
(Theorem~\ref{thm:domain-restriction}).

The fourth paper, \emph{The Geometry of Witnesses}, will explain
why admissibility persists over historical time, introducing
second-order repair operators (witnesses that maintain repair
operators rather than distinctions directly).

The fifth paper, \emph{Conservation of Ambiguity}, will describe
how finite repair resources are allocated across a living distinction
space, with ambiguity, repair cost, and interference as competing
claims on the same budget.

The present paper is designed to serve as a conceptual entry point
to the program.
Its central theorem is stated and proved using only the definitions
in Section~\ref{sec:background} and the Appendix; a reader
encountering the program for the first time may read this paper
first and return to the companion papers for full derivations.

Standard epistemology begins with truth and asks which beliefs
satisfy it.
This ordering is standard but not mandatory.
The repair-theoretic account reverses it: an observer must first
establish which distinctions it can maintain (admissibility), then
identify which of those distinctions it can discriminate in context
(the observable subset), and only then assign truth values.
The ordering $\mathcal{A}(\mathcal{R}) \to \mathcal{D}_O \to
\mathcal{T}|_{\mathcal{D}_O}$ is not a philosophical preference.
It is a theorem.

The paper proceeds as follows.
Section~\ref{sec:background} reviews the repair framework.
Section~\ref{sec:truth} defines truth predicates formally.
Section~\ref{sec:domain} states and proves the Domain Restriction
Theorem.
Section~\ref{sec:sahc} demonstrates that a state-of-the-art
hierarchical classification architecture (SAHC) independently
implements the same structure, and derives the Hierarchy
Consistency Corollary.
Section~\ref{sec:threeval} develops the three-valued semantics.
Sections~\ref{sec:bayes}--\ref{sec:kuhn} derive consequences for
Bayesian inference, machine learning, legal admissibility,
language acquisition, and scientific revolutions.
Section~\ref{sec:limits} shows the ordering does not reverse.
Section~\ref{sec:open} states open problems.
Section~\ref{sec:consequences} draws broader consequences.
The Appendix collects required definitions from the companion
papers.


% ============================================================
\section{Background: Repair, Distinctions, and Admissibility}
\label{sec:background}
% ============================================================

We recall the definitions required from the companion papers.
Full derivations appear in \cite{Flyxion2026a}; we state here only
what is needed for the Domain Restriction Theorem.

A \emph{substrate} $S = (\Xsub, \mathcal{P}, \Dsub)$ consists of
a locally compact Hausdorff state space $\Xsub$, a perturbation
semigroup $\mathcal{P}$, and a closed distinction relation
$\Dsub \subseteq \Xsub \times \Xsub$.
A \emph{repair repertoire} $\mathcal{R} \subseteq C(\Xsub, \Xsub)$
is a sequentially compact set of continuous maps.
The \emph{minimal repair cost} of distinction $d$ under $\mathcal{R}$
is
\[
  \CR(d, \mathcal{R})
  \;=\;
  \sup_{p \in \mathcal{P}_d}
  \inf_{\gamma \in \Gamma_d^p(\mathcal{R})}
  \mathrm{cost}(\gamma),
\]
setting $\CR(d, \mathcal{R}) = +\infty$ when no finite repair path
exists.
The \emph{restorable distinction space} is
\[
  \RestR
  \;=\;
  \bigl\{ d \in \Dsub \;\big|\; \CR(d, \mathcal{R}) < +\infty \bigr\}.
\]
The \emph{admissibility manifold} is
\[
  \AO(\mathcal{R})
  \;=\;
  \overline{\bigl\{
    x \in \Xsub \;\big|\;
    \forall\, d \in \RestR,\;
    \CR(d, \mathcal{R}) < +\infty \text{ from } x
  \bigr\}}.
\]
The Repair Boundary Theorem \cite{Flyxion2026a} states that
$x \in \bA(\mathcal{R})$ if and only if $\CR(d, \mathcal{R}) \to
+\infty$ for some $d \in \RestR$ as the system approaches $x$.

\begin{remark}
  An observer is, in the repair-theoretic account, a pair
  $(S, \mathcal{R})$.
  What the observer can observe is exactly $\RestR$.
  What the observer cannot observe is not hidden from it; it is
  simply not part of the observer's world.
  The admissibility manifold $\AO(\mathcal{R})$ is the geometric
  structure of that world.
\end{remark}


% ============================================================
\section{Truth as a Partial Predicate}
\label{sec:truth}
% ============================================================

The classical notion of a truth predicate takes the form
$\Tsub: \Dsub \to \{0, 1\}$, assigning a truth value to every
distinction in the full distinction space.
We argue that this is the wrong domain.

\begin{definition}[Truth predicate]
  A \emph{truth predicate} for observer $(S, \mathcal{R})$ is a
  function $\Tsub: X \to \{0, 1, \Undef\}$ defined on some subset
  $X \subseteq \Dsub$, where $0$ denotes falsity, $1$ denotes
  truth, and $\Undef$ denotes unevaluability.
  The predicate is \emph{total} if $X = \Dsub$; it is
  \emph{partial} if $X \subsetneq \Dsub$.
\end{definition}

\begin{definition}[Evaluability]
  A distinction $d \in \Dsub$ is \emph{evaluable} by observer
  $(S, \mathcal{R})$ if assigning a truth value to $d$ is
  well-defined, i.e., if $d \in \operatorname{dom}(\Tsub)$ with
  $\Tsub(d) \in \{0,1\}$.
  A distinction is \emph{unevaluable} if $\Tsub(d) = \Undef$.
\end{definition}

The classical picture assumes total truth predicates.
We show this assumption is incoherent for finite observers.

\begin{lemma}[Truth-evaluation is repair]
  Any procedure for evaluating the truth of distinction $d$ must
  be capable of discriminating $d$ from its negation after
  arbitrary admissible perturbation.
  Formally, truth-evaluation of $d$ requires the existence of
  a finite-cost repair path restoring $d$ after perturbation.
  Therefore truth-evaluation of $d$ presupposes
  $d \in \RestR$.
  \label{lem:truth-is-repair}
\end{lemma}

\begin{proof}
  Evaluating whether $d$ holds requires the evaluator to be in a
  state from which $d$ is discriminable: the evaluator must be
  able to distinguish the case $(x,y) \in \Dsub$ from the case
  $(x,y) \notin \Dsub$.
  If $d$ has been erased by an admissible perturbation $p$, the
  evaluator must restore discriminability before it can evaluate.
  The restoration of discriminability after perturbation is
  precisely the repair of $d$: a finite-cost repair path
  $\gamma \in \Gamma_d^p(\mathcal{R})$ returning the system to
  a state where $d$ is discriminable.
  If no such path exists --- if $\CR(d, \mathcal{R}) = +\infty$
  --- then no evaluation procedure can operate, because there is
  no state the evaluator can reach from which the discrimination
  can be made.
  Hence $d \notin \RestR$ implies $d$ is unevaluable.
\end{proof}

\begin{remark}
  Lemma~\ref{lem:truth-is-repair} does not claim that
  truth-evaluation literally requires executing a repair sequence.
  It claims that the \emph{capacity} to evaluate truth requires
  the \emph{existence} of such a sequence.
  An observer that cannot in principle restore $d$ cannot in
  principle evaluate $d$, because it cannot reliably distinguish
  the truth condition from states reached by perturbation.
  This is an operational constraint, not a claim about the
  phenomenology of evaluation.
\end{remark}


% ============================================================
\section{The Domain Restriction Theorem}
\label{sec:domain}
% ============================================================

\begin{theorem}[Domain Restriction Theorem]
  Let $(S, \mathcal{R})$ be an observer and let $\Tsub$ be a
  well-defined truth predicate for this observer.
  Then $\Tsub$ factors through the inclusion
  $\iotamap: \RestR \hookrightarrow \Dsub$:
  \[
    \Tsub \;=\; \Ttilde \circ \iotamap,
  \]
  where $\Ttilde: \RestR \to \{0, 1\}$ is the restriction of
  $\Tsub$ to the restorable distinction space.
  Consequently,
  \[
    \operatorname{dom}(\Tsub) \;\subseteq\; \RestR.
  \]
  A well-defined truth predicate cannot assign truth values outside
  the admissible distinction space.
  \label{thm:domain-restriction}
\end{theorem}

\begin{proof}
  By Lemma~\ref{lem:truth-is-repair}, every distinction in
  $\operatorname{dom}(\Tsub)$ must be in $\RestR$.
  Therefore $\operatorname{dom}(\Tsub) \subseteq \RestR$.

  Define $\Ttilde: \RestR \to \{0,1\}$ by
  $\Ttilde(d) = \Tsub(\iotamap(d)) = \Tsub(d)$ for all
  $d \in \RestR$.
  Since $\iotamap$ is the inclusion map (the identity on its domain),
  $\Tsub = \Ttilde \circ \iotamap$ on $\operatorname{dom}(\Tsub)$.
  For $d \notin \RestR$, $\Tsub(d) = \Undef$ by definition of
  unevaluability, which is consistent with $\Ttilde$ being
  undefined outside $\RestR$.
\end{proof}

\begin{remark}[The categorical reading]
  The factorization $\Tsub = \Ttilde \circ \iotamap$ is a universal
  property: every truth predicate factors through $\iotamap$.
  In categorical terms, $\iotamap: \RestR \hookrightarrow \Dsub$ is
  the unique morphism through which all well-defined truth
  predicates factor.
  This means that the inclusion of the admissible distinction space
  into the full distinction space is the \emph{domain object} for
  truth predicates: any construction involving truth automatically
  inherits the admissibility constraint via this factorization.
\end{remark}

\begin{corollary}[The strict ordering]
  The ordering
  \[
    \AO(\mathcal{R})
    \;\longrightarrow\;
    \RestR
    \;\longrightarrow\;
    \Tsub|_{\RestR}
  \]
  is strict at each arrow.
  Admissibility determines the observer's world;
  restorability determines the domain of truth;
  truth is a predicate on that domain.
  No step in this sequence can be reversed or collapsed without
  contradiction.
  \label{cor:strict-ordering}
\end{corollary}

\begin{proof}
  The first arrow is strict because $\AO(\mathcal{R})$ is a subset
  of $\Xsub$, not of $\Dsub$; the map sends the state manifold to
  the distinction space via the induced discrimination structure.
  The second arrow is strict because $\RestR \subsetneq \Dsub$
  in general: there exist distinctions in $\Dsub$ that are not
  restorable.
  The third arrow is strict because $\Tsub$ is a partial function
  on $\RestR$, not all of $\RestR$ is evaluable in practice even
  if it is admissible in principle.
\end{proof}

\begin{theorem}[Geometry of the evaluability boundary]
  The boundary of $\operatorname{dom}(\Tsub)$ within $\Dsub$
  coincides with the irrestorability boundary: the set of
  distinctions $d$ for which $\CR(d, \mathcal{R}) \to +\infty$.
  The domain of truth therefore inherits the geometric structure
  of the admissibility manifold: it has a well-defined metric
  (the repair pseudometric), a boundary (where repair cost
  diverges), and topological invariants (the homology groups of
  the admissibility manifold).
  \label{thm:evaluability-geometry}
\end{theorem}

\begin{proof}
  $\operatorname{dom}(\Tsub) \subseteq \RestR$ by the Domain
  Restriction Theorem.
  The boundary of $\RestR$ within $\Dsub$ is exactly the set of
  distinctions approaching which $\CR(d, \mathcal{R}) \to +\infty$,
  by the Repair Boundary Theorem of \cite{Flyxion2026a}.
  The repair pseudometric $d_R$ on $\RestR$ and the homology groups
  $H_*(\AO(\mathcal{R}))$ are defined in \cite{Flyxion2026a,Flyxion2026b}
  and are inherited by $\operatorname{dom}(\Tsub)$ via the inclusion
  $\operatorname{dom}(\Tsub) \subseteq \RestR$.
\end{proof}

\begin{remark}
  Theorem~\ref{thm:evaluability-geometry} means that evaluability
  is not a flat binary partition of all possible claims into
  evaluable and unevaluable.
  The domain of truth has boundary structure: some distinctions are
  marginally evaluable (near the irrestorability boundary, with high
  repair cost), and the geometry of that margin is non-trivial.
  Near-boundary evaluability corresponds to claims that can barely
  be assessed, that require enormous epistemic effort, or that
  become unevaluable under small perturbation.
  This is the formal correlate of the philosophical notion of
  ``borderline cases.''
\end{remark}


% ============================================================
\section{Hierarchical Domain Restriction: Empirical Instantiation}
\label{sec:sahc}
% ============================================================

The Domain Restriction Theorem is a mathematical result.
This section argues that it is also an empirical one: a
state-of-the-art deep learning architecture for satellite image
classification independently converged on the same structure,
and its performance gains over non-hierarchical baselines
constitute evidence that the domain restriction is a
genuine design principle, not a philosophical preference.

\subsection{The Architecture as Factorization}

Weikmann et al.~\cite{Weikmann2026} propose the
Semantics-Aware Hierarchical Consensus (SAHC) method
for multi-level remote sensing classification.
Their architecture introduces $H$ classification heads
$g^h \circ f$, one per hierarchical level, together with
trainable hierarchy log-joint matrices $J^{h,h'}$ that project
predictions between levels, and a consensus loss $L_{HC}$
penalising disagreement among all projected predictions at
a target level.
The headline result is consistent improvement over
non-hierarchical baselines across every level of a
90-class taxonomy~\cite{Weikmann2026}.

We claim that this architecture is an empirical instance of
the Domain Restriction Theorem applied at multiple levels
simultaneously.
Let $h = 1$ (coarse), $h = 2$ (intermediate), $h = 3$ (fine)
denote the three classification granularities.
Define:
\begin{itemize}
  \item $\mathcal{R}_1$: the repair repertoire sufficient to
    restore coarse-level distinctions (vegetation vs.\
    non-vegetation; urban vs.\ rural);
  \item $\mathcal{R}_2 \supseteq \mathcal{R}_1$: the repair
    repertoire sufficient to restore intermediate distinctions
    (forest vs.\ shrub; residential vs.\ industrial);
  \item $\mathcal{R}_3 \supseteq \mathcal{R}_2$: the repair
    repertoire sufficient to restore fine-grained distinctions
    (deciduous forest vs.\ coniferous; compact residential
    vs.\ sparse residential).
\end{itemize}
By Repair Monotonicity~\cite{Flyxion2026a},
$\Rest(\mathcal{R}_1) \subseteq \Rest(\mathcal{R}_2) \subseteq
\Rest(\mathcal{R}_3)$,
so the admissibility manifolds are nested:
$\mathcal{A}(\mathcal{R}_1) \subseteq \mathcal{A}(\mathcal{R}_2)
\subseteq \mathcal{A}(\mathcal{R}_3)$.
The Domain Restriction Theorem then yields a chain of
inclusions for the truth predicates at each level:
\[
  \mathcal{T}_f
  \;=\;
  \widetilde{\mathcal{T}}_f
  \circ \iotamap_{f,m}
  \circ \iotamap_{m,c},
\]
where $\iotamap_{m,c}: \Rest(\mathcal{R}_2) \hookrightarrow
\Rest(\mathcal{R}_1)$ and $\iotamap_{f,m}: \Rest(\mathcal{R}_3)
\hookrightarrow \Rest(\mathcal{R}_2)$ are the admissibility-preserving
inclusions.
This says: the fine-level truth predicate $\mathcal{T}_f$
factors through the intermediate admissibility domain,
which in turn factors through the coarse admissibility domain.
A pixel cannot be truthfully classified as ``deciduous forest''
unless it is first admissible as ``forest'' and before that as
``vegetation.''

\subsection{Hierarchy Matrices as Approximate Inclusion Maps}

SAHC's trainable matrices $J^{h,h'} \in [0,1]^{|\Omega^h|
\times |\Omega^{h'}|}$ are initialised from a user-defined binary
hierarchy and then learned from data.
In the repair-theoretic reading, $J^{h,h'}$ is an empirical
approximation of $\iotamap_{h,h'}$: it encodes the learned
conditional probability that a fine-level class $\omega^h_i$
lies within the admissibility domain of a coarser class
$\omega^{h'}_j$.
Where the user-defined hierarchy is rigid, $J^{h,h'}$ learns
which semantic overlaps the rigid taxonomy missed.

The heatmap of the learned joint probability matrix
$J^{h_2,h_1}$ on the ELU dataset (reported in \cite{Weikmann2026})
confirms this interpretation: the highest-weight entries align
with the predefined hierarchy, but off-diagonal weights reveal
genuine semantic adjacencies --- ``water basins'' and
``aquaculture in the marine environment,'' ``continuous urban
fabric'' and ``road networks'' --- that the rigid taxonomy
excluded but the data supports.
These off-diagonal weights are not noise.
They are the learned repair cost relationships: two classes
are semantically adjacent precisely when the repair operators
for one partially restore the other.

\subsection{The Consensus Loss as Domain Restriction Penalty}

SAHC's hierarchical consensus loss $L_{HC}$ penalises the
Jensen--Shannon divergence between the consensus prediction
$\hat{p}^{\tilde{h}}_{HC}$ and the projected predictions from
every other level:
\[
  L_{HC}
  \;=\;
  \frac{1}{N} \sum_n \sum_{\tilde{h}} \frac{1}{\log|\Omega^{\tilde{h}}|}
  \sum_h \mathrm{JSD}^{h,\tilde{h}}(\hat{p}^{\tilde{h}}_{HC} \,\|\,
  \hat{p}^{h \to \tilde{h}}).
\]
In the repair-theoretic reading, this is the penalty for
violating the Domain Restriction Theorem:
\[
  L_{HC}
  \;\approx\;
  d\!\left(
    \widetilde{\mathcal{T}} \circ \iotamap,\;
    \mathcal{T}
  \right),
\]
where $d$ is a divergence between the admissibility-restricted
prediction and the unconstrained prediction.
When $L_{HC} = 0$, the classifier's fine-level outputs are
exactly consistent with their coarse-level projections:
fine distinctions are admissible at every level above them.
When $L_{HC} > 0$, the classifier has assigned a truth value
to a fine distinction that cannot be projected into the coarser
admissibility domains --- it has predicted ``deciduous forest''
where the coarse head predicts ``urban area.''
SAHC penalises this prediction as inadmissible.

The ablation study in \cite{Weikmann2026} isolates this
penalty's contribution: adding $L_{HC}$ to the multi-head
baseline raises fine-grained OA from 96.44\% to 96.52\%,
and adding both consensus training and consensus inference
raises it further to 96.63\%.
The performance gains are not large in absolute terms, but
their mechanism is precisely the domain restriction: forcing
fine-level truth assignments to be restorable through coarser
admissibility levels reduces inadmissible predictions and
improves accuracy.

\subsection{Generalisation: From SAHC to TARTAN}

SAHC assumes a fixed semantic hierarchy: levels are
defined by a human-specified taxonomy, and hierarchy
matrices encode transitions between taxonomy levels.
This is one concrete instance of a more general structure.

In the TARTAN framework, reality is represented as a
hierarchy of interacting tiles at arbitrary resolutions.
Each level carries its own admissibility constraints, and
cross-scale consistency is maintained through transport
operators $T_{i \to j}$ that need not follow a fixed semantic
taxonomy.
The TARTAN coherence functional,
\[
  L_{\mathrm{TARTAN}}
  \;=\;
  \sum_{i < j}
  d\!\left(
    T_{i \to j}(x_i),\; x_j
  \right),
\]
enforces the Domain Restriction Theorem simultaneously
across spatial, temporal, semantic, and computational scales.
SAHC's $J^{h,h'}$ matrices are one specific parameterisation
of $T_{i \to j}$ in the semantic hierarchy case;
SAHC's $L_{HC}$ is one specific instance of $L_{\mathrm{TARTAN}}$
under Jensen--Shannon divergence.

The formal correspondence is:

\begin{center}
\renewcommand{\arraystretch}{1.4}
\begin{tabular}{l@{\quad$\longleftrightarrow$\quad}l}
  Hierarchy levels $h \in [1, H]$ & Resolution levels $i \in I$ \\
  Hierarchy matrix $J^{h,h'}$ & Transport operator $T_{i \to j}$ \\
  Projected prediction $\hat{p}^{h \to h'}$ & Transported state $T_{i \to j}(x_i)$ \\
  Consensus loss $L_{HC}$ & Coherence functional $L_{\mathrm{TARTAN}}$ \\
  Semantic taxonomy & Arbitrary resolution hierarchy \\
\end{tabular}
\end{center}

SAHC is therefore not merely a hierarchical classifier
but a specific instance of admissibility-preserving multiscale
computation.
Its empirical success across two datasets with very different
hierarchical complexity suggests that the Domain Restriction
Theorem is not only a logical consequence of the repair
framework but a productive engineering principle when
implemented at scale.

\subsection{A Formal Consequence}

\begin{corollary}[Hierarchy Consistency Corollary]
  Let $\mathcal{R}_1 \subseteq \mathcal{R}_2 \subseteq \cdots
  \subseteq \mathcal{R}_H$ be a hierarchy of repair repertoires
  with admissibility-preserving inclusions
  $\iotamap_{h,h'}: \Rest(\mathcal{R}_h) \hookrightarrow
  \Rest(\mathcal{R}_{h'})$ for all $h > h'$.
  Suppose a prediction architecture produces truth predicates
  $\mathcal{T}_h$ at each level satisfying
  \[
    \mathcal{T}_h \;=\; \widetilde{\mathcal{T}}_h \circ \iotamap_{h,h'}
    \quad \text{for all } h > h'.
  \]
  Then each $\mathcal{T}_h$ satisfies the Domain Restriction Theorem
  at its own level, and the composition
  $\mathcal{T}_H = \widetilde{\mathcal{T}}_H \circ \iotamap_{H,H-1}
  \circ \cdots \circ \iotamap_{2,1}$
  satisfies it simultaneously at all levels.
  \label{cor:hierarchy-consistency}
\end{corollary}

\begin{proof}
  By the Domain Restriction Theorem applied to each level:
  $\operatorname{dom}(\mathcal{T}_h) \subseteq \Rest(\mathcal{R}_h)$
  for each $h$.
  Since $\mathcal{T}_h = \widetilde{\mathcal{T}}_h \circ \iotamap_{h,h'}$
  and $\iotamap_{h,h'}$ maps $\Rest(\mathcal{R}_h)$ into
  $\Rest(\mathcal{R}_{h'})$, evaluability at level $h$ implies
  evaluability at level $h'$.
  The composed factorization follows by induction on the
  number of levels.
\end{proof}

\begin{remark}
  Corollary~\ref{cor:hierarchy-consistency} is the formal
  statement that SAHC-style architectures, when their hierarchy
  matrices exactly implement the admissibility-preserving
  inclusions $\iotamap_{h,h'}$, simultaneously satisfy the
  Domain Restriction Theorem at every level.
  SAHC's learned matrices $J^{h,h'}$ are approximate
  implementations of $\iotamap_{h,h'}$; the consensus loss
  $L_{HC}$ is the gradient signal that drives them toward
  exactness.
  The residual violation --- $L_{HC} > 0$ at convergence
  --- measures the extent to which the architecture departs
  from strict multi-level domain restriction.
  Architectures with smaller $L_{HC}$ at convergence are, in
  this precise sense, better implementations of the theorem.
\end{remark}



The Domain Restriction Theorem introduces $\Undef$ as a third
truth value for unevaluable distinctions.
This is not a mere convenience; the three-valued structure follows
necessarily from the repair framework.

\begin{proposition}[Unevaluability is not falsity]
  If $d \notin \RestR$, then $\Tsub(d) \neq 0$.
  An unevaluable distinction is not false.
  \label{prop:bot-not-false}
\end{proposition}

\begin{proof}
  Assigning $\Tsub(d) = 0$ (falsity) would itself be a truth
  assignment, which by the Domain Restriction Theorem requires
  $d \in \operatorname{dom}(\Tsub) \subseteq \RestR$.
  But we assumed $d \notin \RestR$.
  Contradiction.
  Therefore $\Tsub(d) \neq 0$.
  By the same argument, $\Tsub(d) \neq 1$.
  Hence $\Tsub(d) = \Undef$.
\end{proof}

\begin{remark}
  Proposition~\ref{prop:bot-not-false} resolves a persistent
  confusion in accounts of semantic unevaluability.
  Unevaluable distinctions are sometimes treated as ``false by
  default'' (no truth-maker, no truth) or as ``indeterminate''
  (neither true nor false).
  The repair-theoretic account gives a precise meaning to neither:
  $\Undef$ is not a weakened form of falsity or truth but a distinct
  value arising from the failure of the precondition for evaluation.
  The natural semantic framework is Kleene's strong three-valued
  logic~\cite{Kleene1952}, in which $\Undef$ propagates through
  connectives in a specific way.
  Whether the repair-theoretic $\Undef$ satisfies the Kleene truth
  tables exactly is a question for the logic paper this framework
  suggests; we leave it as an open problem.
\end{remark}

\begin{proposition}[Admissibility expansion increases evaluability]
  If $\mathcal{R}_1 \subseteq \mathcal{R}_2$, then
  $\operatorname{dom}(\Tsub^{(1)}) \subseteq
  \operatorname{dom}(\Tsub^{(2)})$.
  Acquiring more repair operators strictly expands the domain
  of truth.
  \label{prop:expansion}
\end{proposition}

\begin{proof}
  By the Repair Monotonicity Theorem of \cite{Flyxion2026a},
  $\mathcal{R}_1 \subseteq \mathcal{R}_2$ implies
  $\Rest(\mathcal{R}_1) \subseteq \Rest(\mathcal{R}_2)$.
  Since $\operatorname{dom}(\Tsub^{(i)}) \subseteq
  \Rest(\mathcal{R}_i)$, the inclusion follows.
\end{proof}

\begin{remark}
  Proposition~\ref{prop:expansion} is the logical analogue of
  the ``Learning enlarges worlds'' corollary from
  \cite{Flyxion2026a}.
  There it meant that acquiring repair operators expands the
  observer's perceptual world.
  Here it means acquiring repair operators expands the domain of
  truth: claims that were unevaluable become evaluable.
  Learning is not merely acquiring new true beliefs; it is
  expanding the space of evaluable claims.
\end{remark}


% ============================================================
\section{Bayesian Inference}
\label{sec:bayes}
% ============================================================

Standard Bayesian inference assigns a prior $P: \Dsub \to [0,1]$
to every distinction in the full distinction space, with
$P(d)$ interpreted as the probability that $d$ is true.
The Domain Restriction Theorem imposes a constraint on this picture
that has not previously been made explicit.

\begin{proposition}[Coherence requires admissibility]
  A Bayesian prior $P$ is coherent for observer $(S, \mathcal{R})$
  only if $\operatorname{supp}(P) \subseteq \RestR$.
  Assigning a non-trivial prior to an inadmissible distinction is
  incoherent.
  \label{prop:bayes}
\end{proposition}

\begin{proof}
  A non-trivial prior $P(d) \in (0,1)$ presupposes that $d$ can
  in principle be evaluated: otherwise the probability cannot be
  updated by evidence.
  But by the Domain Restriction Theorem, $d$ is evaluable only if
  $d \in \RestR$.
  Therefore $P(d) \notin \{0, 1\}$ requires $d \in \RestR$, and
  the support of any coherent prior lies within $\RestR$.
\end{proof}

\begin{remark}
  Proposition~\ref{prop:bayes} does not prohibit assigning $P(d)=0$
  to inadmissible distinctions; it prohibits treating that
  assignment as meaningful.
  A prior that places zero probability on $d \notin \RestR$ is
  technically coherent but trivially so: it says nothing about $d$
  because $d$ cannot be evaluated.
  The operationally meaningful content of a Bayesian prior is
  entirely concentrated on $\RestR$.

  This has a practical consequence: bounded rationality is not a
  cognitive limitation requiring separate modelling.
  It is the necessary consequence of finite $\RestR$.
  A perfectly rational Bayesian agent with finite repair capacity
  has a prior supported on a proper subset of all logically possible
  distinctions.
  The shape of that subset is determined by the agent's admissibility
  manifold, not by computational constraints.
\end{remark}


% ============================================================
\section{Machine Learning}
\label{sec:ml}
% ============================================================

A classifier trained to distinguish class $A$ from class $B$ is,
in repair-theoretic terms, constructing an approximation to a repair
operator for the distinction $d = (A, B)$.
Training is the process of building $\mathcal{R}$ such that
$d \in \RestR$.

\begin{proposition}[Out-of-distribution failure is inadmissibility]
  Let $d_{\text{train}}$ be the distinction for which a classifier
  was trained and let $d_{\text{deploy}}$ be the distinction it
  faces at deployment.
  Out-of-distribution failure occurs when
  $d_{\text{deploy}} \notin \Rest(\mathcal{R}_{\text{train}})$:
  the distinction to be evaluated was not in the admissibility
  manifold of the training repair repertoire.
  \label{prop:ood}
\end{proposition}

\begin{proof}
  Training builds repair operators sufficient to restore
  $d_{\text{train}}$ after the perturbations seen during training.
  At deployment, if $d_{\text{deploy}}$ differs from
  $d_{\text{train}}$ in ways that require repair operators not
  in $\mathcal{R}_{\text{train}}$, then $\CR(d_{\text{deploy}},
  \mathcal{R}_{\text{train}}) = +\infty$ and the classifier
  cannot evaluate $d_{\text{deploy}}$.
  This is the formal content of out-of-distribution failure.
\end{proof}

\begin{remark}
  The standard account of distribution shift appeals to statistical
  distance between training and deployment distributions.
  The repair-theoretic account gives a geometric alternative: the
  relevant quantity is whether the deployment distinction lies
  within or outside the admissibility manifold of the trained
  repair repertoire.
  These accounts are not equivalent.
  Two distributions can be statistically close while the deployment
  distinction lies outside the training admissibility manifold
  (if the distributions agree on most features but differ on the
  specific repair-relevant features).
  Conversely, distributions can be statistically distant while the
  deployment distinction remains admissible.
  The geometric account is more fundamental.
\end{remark}


% ============================================================
\section{Legal Admissibility}
\label{sec:legal}
% ============================================================

Legal systems provide the most explicit real-world instance of the
domain restriction structure.
Rules of evidence define which distinctions are admissible for
truth-evaluation in legal proceedings.
The phrase ``admissibility of evidence'' is not a metaphor for the
formal structure developed here; it is the same structure at the
institutional level.

\begin{definition}[Legal admissibility manifold]
  The \emph{legal admissibility manifold} of a legal system is
  $\AO(\mathcal{R}_{\text{legal}})$, where
  $\mathcal{R}_{\text{legal}}$ consists of the evidentiary and
  procedural rules that constitute the system's repair repertoire.
  A distinction $d$ between legal and illegal conduct is evaluable
  by the court if and only if $d \in \Rest(\mathcal{R}_{\text{legal}})$.
\end{definition}

\begin{proposition}[Verdicts are domain-restricted]
  A legal verdict assigns a truth value ($1$ = guilty,
  $0$ = not guilty) to a distinction $d$.
  By the Domain Restriction Theorem, this is only coherent if
  $d \in \Rest(\mathcal{R}_{\text{legal}})$.
  Evidence that falls outside the legal admissibility manifold
  does not make a verdict false; it makes the verdict unevaluable
  through that evidence.
\end{proposition}

\begin{remark}
  This accounts for two features of legal reasoning that seem
  paradoxical from a purely truth-theoretic perspective.
  First, an inadmissible piece of evidence is not excluded because
  it is false or unreliable; it is excluded because the legal
  system's repair repertoire cannot incorporate it.
  The distinction it would resolve is not within
  $\Rest(\mathcal{R}_{\text{legal}})$.
  Second, a verdict of ``not guilty'' is not a verdict of
  ``innocent''; it is a verdict that the distinction between guilt
  and innocence was not established within the admissibility
  manifold of the proceeding.
  Both features follow immediately from the domain restriction.
\end{remark}


% ============================================================
\section{Language Acquisition}
\label{sec:language}
% ============================================================

A child acquiring a language is expanding $\RestR$ in a specific
domain: the phonological, syntactic, and semantic distinctions
that the language encodes.

\begin{proposition}[Grammaticality is admissibility]
  A sentence $s$ is grammatical in language $L$ if and only if
  the distinctions $s$ expresses are in $\Rest(\mathcal{R}_L)$,
  where $\mathcal{R}_L$ is the repair repertoire associated with
  competence in $L$.
  An ungrammatical sentence does not express a false proposition;
  it fails to express a proposition at all, receiving $\Tsub(s)
  = \Undef$.
  \label{prop:grammar}
\end{proposition}

\begin{proof}
  Grammaticality is the condition under which a sentence picks out
  a recoverable distinction in the language's semantic space.
  An ungrammatical sentence is one for which no repair path exists
  from the noisy or degraded signal to a well-formed distinction:
  $\CR(d_s, \mathcal{R}_L) = +\infty$.
  By the Domain Restriction Theorem, $\Tsub(d_s) = \Undef$.
\end{proof}

\begin{remark}
  Proposition~\ref{prop:grammar} makes precise a claim that
  linguists and philosophers have approached informally: that
  truth-conditional semantics presupposes grammaticality.
  A sentence must be grammatical before it can be true or false.
  The repair-theoretic account gives this presupposition a
  formal home: grammaticality is admissibility, and admissibility
  is prior to truth by the domain restriction.

  Language acquisition is repair-repertoire expansion.
  A child who learns a new grammatical construction acquires a
  new class of repair operators, expanding $\RestR$ and therefore
  expanding $\operatorname{dom}(\Tsub)$.
  The child gains not merely new true beliefs but the capacity
  to evaluate a larger class of claims.
\end{remark}


% ============================================================
\section{Kuhn Without Relativism}
\label{sec:kuhn}
% ============================================================

The Kuhnian account of scientific revolutions holds that scientists
in different paradigms inhabit different conceptual worlds and that
their claims are incommensurable.
Critics have argued that this collapses into relativism: if there
is no shared domain of evaluation, there is no fact of the matter
about which paradigm is better.

The repair-theoretic account dissolves this debate by making the
structure precise.

\begin{theorem}[Paradigm shifts as domain transitions]
  Let $\mathcal{R}_{\text{old}}$ and $\mathcal{R}_{\text{new}}$
  be the repair repertoires of two paradigms.
  A non-trivial paradigm shift satisfies
  \[
    \Rest(\mathcal{R}_{\text{old}})
    \;\triangle\;
    \Rest(\mathcal{R}_{\text{new}})
    \;\neq\; \emptyset,
  \]
  where $\triangle$ denotes symmetric difference.
  By the Domain Restriction Theorem, the domains of truth differ:
  \[
    \operatorname{dom}(\Tsub^{\text{old}})
    \;\triangle\;
    \operatorname{dom}(\Tsub^{\text{new}})
    \;\neq\; \emptyset.
  \]
  Some previously evaluable claims become unevaluable;
  some previously unevaluable claims become evaluable.
  \label{thm:paradigm-shift}
\end{theorem}

\begin{proof}
  The first statement is the Paradigm Shift Theorem proved in
  \cite{Flyxion2026a}.
  The second follows immediately by applying the Domain Restriction
  Theorem to both repertoires.
\end{proof}

\begin{remark}
  Theorem~\ref{thm:paradigm-shift} gives precise content to
  Kuhnian incommensurability without relativism.
  Incommensurability is not the claim that two paradigms disagree
  about the same facts.
  It is the structural claim that their domains of truth differ:
  $\operatorname{dom}(\Tsub^{\text{old}}) \triangle
  \operatorname{dom}(\Tsub^{\text{new}}) \neq \emptyset$.
  Claims in the symmetric difference are evaluable in one paradigm
  and unevaluable in the other.

  This is not relativism because the structure is observer-relative,
  not truth-relative.
  There is no claim that the same distinction is both true and false.
  There is only the structural fact that different repair repertoires
  generate different domains of evaluability.
  Whether one paradigm is better than another is a question about
  which admissibility manifold is more productive, more robust,
  or more adequate to the phenomena — a comparison that is
  perfectly meaningful without requiring a view from nowhere.
\end{remark}


% ============================================================
\section{The Limits of the Ordering}
\label{sec:limits}
% ============================================================

The Domain Restriction Theorem establishes that admissibility
precedes truth.
A natural question is whether this ordering can be reversed:
can truth-evaluation expand the admissibility manifold?

\begin{proposition}[The ordering does not reverse]
  Truth-evaluation cannot expand $\RestR$.
  Evaluating the truth of $d \in \RestR$ does not add any
  distinction to $\RestR$.
  \label{prop:no-reversal}
\end{proposition}

\begin{proof}
  Truth-evaluation operates on distinctions already in
  $\RestR$: it assigns $0$ or $1$ to $d$.
  It does not create new repair operators.
  Therefore $\RestR$ is unchanged by truth-evaluation.
  Only acquiring new repair operators (expanding $\mathcal{R}$)
  can expand $\RestR$, and that is a separate act from evaluation.
\end{proof}

\begin{remark}
  The ordering $\mathcal{A}(\mathcal{R}) \to \RestR \to
  \Tsub|_{\RestR}$ therefore forms a strict hierarchy with no
  feedback from truth to admissibility.
  Feedback does exist in practice: discovering that a distinction
  is evaluable and important can motivate acquiring new repair
  operators, thereby expanding admissibility.
  But this feedback passes through the agent's actions, not through
  the logical structure of truth-evaluation itself.
  The logical ordering is irreversible; the causal story is richer.
\end{remark}


% ============================================================
\section{Open Problems}
\label{sec:open}
% ============================================================

\begin{enumerate}

  \item \textbf{Three-valued semantics.}\;
    Does the repair-theoretic $\Undef$ satisfy Kleene's strong
    three-valued truth tables exactly, or does it require a
    different connective structure?
    The propagation of $\Undef$ through conjunction, disjunction,
    and negation in the repair framework may differ from Kleene's
    tables in cases where one conjunct is evaluable and the other
    is not.
    A formal comparison with paraconsistent and partial-logic
    frameworks~\cite{Kleene1952} is needed.

  \item \textbf{The topology of evaluability.}\;
    Theorem~\ref{thm:evaluability-geometry} establishes that
    $\operatorname{dom}(\Tsub)$ inherits the topology of
    $\AO(\mathcal{R})$.
    What are the topological invariants of the evaluability domain
    for various natural observers (logical agents, scientific
    communities, legal systems)?
    Is the evaluability domain always connected?
    What is the significance of non-trivial $H_1(\operatorname{dom}(\Tsub))$?

  \item \textbf{Degrees of evaluability.}\;
    The current framework is binary: a distinction is either
    evaluable or not.
    A graded generalisation would define evaluability as a function
    of $\CR(d, \mathcal{R})$: high repair cost corresponds to
    marginal evaluability, low repair cost to robust evaluability.
    This would give a natural account of ``hard cases'' in law,
    ``borderline cases'' in vagueness, and ``underdetermined''
    hypotheses in philosophy of science.

  \item \textbf{Joint evaluability and interference.}\;
    When is a set of distinctions jointly evaluable?
    The interference matrix from \cite{Flyxion2026a} measures
    repair conflict between distinctions.
    High interference ($I(d_1, d_2) > 0$) means jointly evaluating
    $d_1$ and $d_2$ costs more than evaluating each separately.
    Is there a notion of a \emph{consistent evaluation domain}:
    a maximal set of distinctions that can be jointly evaluated
    within a given repair budget?

  \item \textbf{Truth predicates across paradigms.}\;
    Theorem~\ref{thm:paradigm-shift} identifies the symmetric
    difference of evaluability domains across a paradigm shift.
    Can the repair pseudometric be used to define a distance between
    paradigms: how much does the evaluability domain change?
    A small symmetric difference with low repair cost corresponds
    to a minor revision; a large or topologically complex one
    corresponds to a major revolution.

\end{enumerate}


% ============================================================
\section{Consequences}
\label{sec:consequences}
% ============================================================

Taking the Domain Restriction Theorem seriously reorganises a
range of standing debates.

\medskip\noindent
\textbf{Epistemology.}\;
The classical problem of the criterion --- how do we know what we
know without already knowing how to know it --- is dissolved by
the domain restriction.
The criterion is not epistemic but operational: the admissibility
manifold $\AO(\mathcal{R})$ is determined by the observer's repair
capacity, not by a prior epistemic standard.
The regress terminates not in a foundational truth but in a repair
repertoire.

\medskip\noindent
\textbf{Philosophy of science.}\;
Scientific progress is not merely the accumulation of true beliefs.
It is the expansion of the domain of evaluability: acquiring repair
operators that make previously unevaluable distinctions evaluable.
The history of science is the history of $\RestR$ growing.

\medskip\noindent
\textbf{Philosophy of language.}\;
Meaning is not prior to truth; admissibility is prior to both.
A meaningful sentence is one that picks out an admissible
distinction; a true sentence is one that correctly characterises
that distinction.
The admissibility manifold is the formal home for what Wittgenstein
called the limits of one's world.

\medskip\noindent
\textbf{Formal logic.}\;
Classical two-valued logic is the limiting case in which the
observer's repair repertoire is unlimited: $\RestR = \Dsub$.
Three-valued and paraconsistent logics are not alternatives to
classical logic but natural consequences of finite repair capacity.
The logical framework appropriate to a finite observer is
determined by its admissibility manifold.

\medskip
Admissibility is not a property truth has or lacks.
It is the precondition for truth having anything to say.


% ============================================================
\appendix
\section{Summary of Required Definitions}
\label{app:defs}
% ============================================================

This appendix collects the definitions from \cite{Flyxion2026a}
required for the main results.
It is included so that the Domain Restriction Theorem
(Theorem~\ref{thm:domain-restriction}) can be read independently.
Full derivations, axioms, and proofs appear in \cite{Flyxion2026a}.

\medskip\noindent
\textbf{A.1 Substrate.}\;
A \emph{substrate} $S = (\Xsub, \mathcal{P}, \Dsub)$ consists of
a locally compact Hausdorff state space $\Xsub$, a perturbation
semigroup $\mathcal{P}$, and a closed distinction relation
$\Dsub \subseteq \Xsub \times \Xsub$.

\medskip\noindent
\textbf{A.2 Repair repertoire.}\;
A \emph{repair repertoire} $\mathcal{R} \subseteq C(\Xsub, \Xsub)$
is a sequentially compact set of continuous maps (repair operators)
under the compact-open topology.

\medskip\noindent
\textbf{A.3 Minimal repair cost.}\;
$\CR(d, \mathcal{R}) = \sup_{p \in \mathcal{P}_d}
\inf_{\gamma \in \Gamma_d^p(\mathcal{R})} \mathrm{cost}(\gamma)$,
with $\CR(d, \mathcal{R}) = +\infty$ if no finite repair path
exists.

\medskip\noindent
\textbf{A.4 Restorable distinction space.}\;
$\RestR = \{d \in \Dsub \mid \CR(d, \mathcal{R}) < +\infty\}$.

\medskip\noindent
\textbf{A.5 Admissibility manifold.}\;
$\AO(\mathcal{R}) = \overline{\{x \in \Xsub \mid \forall\,
d \in \RestR,\, \CR(d,\mathcal{R}) < +\infty \text{ from } x\}}$.

\medskip\noindent
\textbf{A.6 Repair pseudometric.}\;
$d_R(d_i, d_j) = \inf_{\gamma: d_i \rightsquigarrow d_j}
\int_\gamma c\, \mathrm{d}s$, defining a metric on $\RestR$ in
the reversible case.

\medskip\noindent
\textbf{A.7 Repair Boundary Theorem} \emph{(proved in \cite{Flyxion2026a})}.
$x \in \bA(\mathcal{R})$ if and only if
$\CR(d, \mathcal{R}) \to +\infty$ for some $d \in \RestR$
as the system approaches $x$.

\medskip\noindent
\textbf{A.8 Repair Monotonicity Theorem}
\emph{(proved in \cite{Flyxion2026a})}.
$\mathcal{R}_1 \subseteq \mathcal{R}_2$ implies
$\Rest(\mathcal{R}_1) \subseteq \Rest(\mathcal{R}_2)$.

\medskip\noindent
\textbf{A.9 Paradigm Shift Theorem}
\emph{(proved in \cite{Flyxion2026a})}.
A non-trivial repair-repertoire change produces
$\Rest(\mathcal{R}_{\text{old}}) \triangle
\Rest(\mathcal{R}_{\text{new}}) \neq \emptyset$.


% ============================================================
\section{The Evaluability Measure}
\label{app:measure}
% ============================================================

The main text treats evaluability as binary: a distinction $d$ is
either evaluable ($d \in \RestR$) or unevaluable ($d \notin \RestR$).
This appendix develops a graded version in which evaluability is a
real-valued measure on the distinction space, vanishing at the
irrestorability boundary and positive in the interior.

\begin{definition}[Evaluability measure]
  Let $f: [0, +\infty] \to [0, 1]$ be a continuous, strictly
  decreasing function with $f(0) = 1$ and $\lim_{c \to \infty}
  f(c) = 0$.
  The \emph{evaluability measure} of distinction $d$ under
  repertoire $\mathcal{R}$ is
  \[
    \mu_{\mathcal{R}}(d)
    \;=\;
    f\!\bigl(\CR(d, \mathcal{R})\bigr)
    \;\in\; [0, 1].
  \]
  Natural choices include $f(c) = e^{-c}$ (exponential decay)
  and $f(c) = (1 + c)^{-1}$ (hyperbolic decay).
\end{definition}

The evaluability measure extends the binary domain restriction to
a soft version: $\mu_{\mathcal{R}}(d) = 1$ for costlessly
restorable distinctions, $\mu_{\mathcal{R}}(d) \approx 0$ for
near-irrestorable distinctions, and $\mu_{\mathcal{R}}(d) = 0$
precisely when $d \notin \RestR$.
The binary Domain Restriction Theorem is the $\{0,1\}$-valued
special case obtained by thresholding at any level $\epsilon > 0$.

\begin{proposition}[Monotonicity of the evaluability measure]
  If $\mathcal{R}_1 \subseteq \mathcal{R}_2$, then
  $\mu_{\mathcal{R}_1}(d) \leq \mu_{\mathcal{R}_2}(d)$
  for all $d \in \Dsub$.
  Acquiring more repair operators increases evaluability pointwise.
  \label{prop:measure-monotone}
\end{proposition}

\begin{proof}
  By Repair Monotonicity~\cite{Flyxion2026a},
  $\CR(d, \mathcal{R}_2) \leq \CR(d, \mathcal{R}_1)$ for all $d$:
  a larger repertoire can only reduce or maintain minimal repair cost.
  Since $f$ is strictly decreasing,
  $f(\CR(d, \mathcal{R}_2)) \geq f(\CR(d, \mathcal{R}_1))$,
  giving $\mu_{\mathcal{R}_2}(d) \geq \mu_{\mathcal{R}_1}(d)$.
\end{proof}

\begin{proposition}[Boundary behaviour]
  $\mu_{\mathcal{R}}(d) \to 0$ as $d$ approaches the
  irrestorability boundary $\partial\AO(\mathcal{R})$ in the
  repair pseudometric.
  \label{prop:measure-boundary}
\end{proposition}

\begin{proof}
  By the Repair Boundary Theorem~\cite{Flyxion2026a},
  $\CR(d, \mathcal{R}) \to +\infty$ as $d \to \partial\AO(\mathcal{R})$.
  Since $f(c) \to 0$ as $c \to +\infty$,
  $\mu_{\mathcal{R}}(d) = f(\CR(d, \mathcal{R})) \to 0$.
\end{proof}

\begin{definition}[Total evaluability]
  Let $\nu$ be a reference measure on $\Dsub$.
  The \emph{total evaluability} of observer $(S, \mathcal{R})$ is
  \[
    E(\mathcal{R})
    \;=\;
    \int_{\Dsub} \mu_{\mathcal{R}}(d)\, \mathrm{d}\nu(d)
    \;=\;
    \int_{\Dsub} f(\CR(d, \mathcal{R}))\, \mathrm{d}\nu(d).
  \]
\end{definition}

\begin{proposition}[Total evaluability and repair-repertoire reduction]
  Let $(\mathcal{R}_t)_{t \geq 0}$ be an erasure path in the sense
  of \cite{Flyxion2026b}.
  Then $E(\mathcal{R}_t)$ is non-increasing in $t$:
  progressive erasure weakly reduces total evaluability.
  \label{prop:total-ev-decreasing}
\end{proposition}

\begin{proof}
  For each $d$, $\CR(d, \mathcal{R}_t)$ is non-decreasing in $t$
  along the erasure path (by the monotonicity of hysteresis under
  progressive erasure, proved in \cite{Flyxion2026b}).
  Since $f$ is decreasing, $\mu_{\mathcal{R}_t}(d)$ is non-increasing
  in $t$ for each $d$.
  By dominated convergence (each integrand is bounded in $[0,1]$),
  $E(\mathcal{R}_t) = \int \mu_{\mathcal{R}_t}\, \mathrm{d}\nu$ is
  non-increasing.
\end{proof}

\begin{remark}
  Proposition~\ref{prop:total-ev-decreasing} connects the evaluability
  measure to the erasure geometry of \cite{Flyxion2026b}: the total
  evaluability $E(\mathcal{R}_t)$ is the integral of the evaluability
  measure over the distinction space, and it decreases monotonically
  as the admissibility manifold deforms under erasure.
  The rate $-\mathrm{d}E/\mathrm{d}t$ is a scalar measure of how fast
  the observer's domain of truth is shrinking --- a global version of
  the boundary migration measure defined in \cite{Flyxion2026b}.
  When $f = e^{-c}$, total evaluability $E(\mathcal{R})$ is the
  Laplace transform of the repair cost distribution, evaluated at $s=1$,
  giving a direct connection to the moment-generating function of $\CR$.
\end{remark}


% ============================================================
\section{Domain Restriction in Continuous Resolution Towers}
\label{app:tartan}
% ============================================================

Section~\ref{sec:sahc} showed that SAHC's finite hierarchy
$h \in \{1, \ldots, H\}$ instantiates the Domain Restriction
Theorem at each level.
This appendix extends the Hierarchy Consistency Corollary
(Corollary~\ref{cor:hierarchy-consistency}) to continuous resolution
towers, which is the setting required for the full TARTAN
framework.

\begin{definition}[Resolution tower]
  A \emph{resolution tower} is a directed family of repair
  repertoires $(\mathcal{R}_s)_{s \in [0,\infty)}$ indexed by
  a continuous resolution parameter $s$, with
  $\mathcal{R}_s \subseteq \mathcal{R}_{s'}$ whenever $s \leq s'$.
  Higher $s$ corresponds to finer resolution.
  The associated restorable distinction spaces form a nested family:
  $\Rest(\mathcal{R}_s) \subseteq \Rest(\mathcal{R}_{s'})$
  for $s \leq s'$, by Repair Monotonicity.
\end{definition}

\begin{definition}[Transport operator]
  A \emph{transport operator} between resolutions $s \leq s'$
  is a map $T_{s \to s'}: \Rest(\mathcal{R}_s) \to
  \Rest(\mathcal{R}_{s'})$ satisfying:
  \begin{enumerate}
    \item $T_{s \to s} = \mathrm{id}$ (identity at equal resolutions);
    \item $T_{s \to s''} = T_{s' \to s''} \circ T_{s \to s'}$
      for all $s \leq s' \leq s''$ (cocycle condition);
    \item $T_{s \to s'}$ is admissibility-preserving: if
      $\Tsub_s(d) \in \{0,1\}$ then $\Tsub_{s'}(T_{s \to s'}(d))
      \in \{0,1\}$.
  \end{enumerate}
\end{definition}

The cocycle condition makes the family $(T_{s \to s'})$ into a
continuous analogue of the composition
$\iotamap_{H,H-1} \circ \cdots \circ \iotamap_{2,1}$
from the finite hierarchy case.
Condition (3) is the continuous version of admissibility-preservation.

\begin{theorem}[Continuous Hierarchy Consistency Theorem]
  Let $(\mathcal{R}_s)_{s \in [0,\infty)}$ be a resolution tower
  with transport operators $(T_{s \to s'})$.
  Suppose a family of truth predicates $(\Tsub_s)_{s \geq 0}$
  satisfies
  \[
    \Tsub_s(d) \;=\; \Tsub_{s'}(T_{s \to s'}(d))
    \quad\text{for all } s \leq s' \text{ and all }
    d \in \Rest(\mathcal{R}_s).
  \]
  Then each $\Tsub_s$ satisfies the Domain Restriction Theorem:
  $\operatorname{dom}(\Tsub_s) \subseteq \Rest(\mathcal{R}_s)$.
  Moreover, the assignment $s \mapsto \operatorname{dom}(\Tsub_s)$
  is monotone: coarser resolutions have smaller evaluability domains.
  \label{thm:continuous-hierarchy}
\end{theorem}

\begin{proof}
  Fix $s \geq 0$.
  Suppose $d \in \operatorname{dom}(\Tsub_s)$, so
  $\Tsub_s(d) \in \{0,1\}$.
  We must show $d \in \Rest(\mathcal{R}_s)$.

  By Lemma~\ref{lem:truth-is-repair} applied at resolution $s$,
  truth-evaluation of $d$ at resolution $s$ requires the existence
  of a finite-cost repair path under $\mathcal{R}_s$.
  Therefore $\CR(d, \mathcal{R}_s) < +\infty$, so
  $d \in \Rest(\mathcal{R}_s)$.

  Monotonicity: for $s \leq s'$, if $d \in \operatorname{dom}(\Tsub_s)$
  then $d \in \Rest(\mathcal{R}_s) \subseteq \Rest(\mathcal{R}_{s'})$,
  so $T_{s \to s'}(d) \in \Rest(\mathcal{R}_{s'})$ and
  $\Tsub_{s'}(T_{s \to s'}(d)) = \Tsub_s(d) \in \{0,1\}$,
  giving $T_{s \to s'}(d) \in \operatorname{dom}(\Tsub_{s'})$.
  Therefore the evaluability domain grows with resolution:
  $T_{s \to s'}(\operatorname{dom}(\Tsub_s)) \subseteq
  \operatorname{dom}(\Tsub_{s'})$.
\end{proof}

\begin{definition}[TARTAN coherence functional]
  Given a resolution tower with transport operators, the
  \emph{TARTAN coherence functional} at resolutions
  $s_1 < s_2 < \cdots < s_n$ is
  \[
    L_{\mathrm{TARTAN}}
    \;=\;
    \sum_{i < j}
    d_R\!\left(
      T_{s_i \to s_j}(x_{s_i}),\; x_{s_j}
    \right),
  \]
  where $d_R$ is the repair pseudometric and $x_s$ denotes the
  observer's distinction-space representation at resolution $s$.
  \label{def:tartan-loss}
\end{definition}

\begin{proposition}[Zero coherence implies consistent domain restriction]
  If $L_{\mathrm{TARTAN}} = 0$ for all pairs $(s_i, s_j)$,
  then the truth predicates $(\Tsub_s)$ satisfy the consistency
  condition of Theorem~\ref{thm:continuous-hierarchy} and the
  Domain Restriction Theorem holds at every resolution.
  \label{prop:tartan-zero}
\end{proposition}

\begin{proof}
  $L_{\mathrm{TARTAN}} = 0$ implies $T_{s_i \to s_j}(x_{s_i}) =
  x_{s_j}$ in the repair pseudometric for all pairs, which is
  exactly the consistency condition $\Tsub_{s_i} = \Tsub_{s_j}
  \circ T_{s_i \to s_j}$.
  Theorem~\ref{thm:continuous-hierarchy} then applies.
\end{proof}

\begin{remark}
  Proposition~\ref{prop:tartan-zero} gives a precise sense in which
  $L_{\mathrm{TARTAN}}$ is a penalty for violating the Domain
  Restriction Theorem across the continuous resolution tower.
  Minimising $L_{\mathrm{TARTAN}}$ during training drives the
  system toward a state where truth assignments at every resolution
  are mutually consistent and each is restricted to its own
  admissibility domain.
  SAHC's consensus loss $L_{HC}$ is the finite, discrete-hierarchy
  instance of $L_{\mathrm{TARTAN}}$: it penalises the same
  violation at finitely many resolution levels using Jensen--Shannon
  divergence as the specific choice of $d_R$.
  The continuous version in Definition~\ref{def:tartan-loss} applies
  equally to spatial resolution towers (pixels at $10\,\mathrm{m}$,
  $100\,\mathrm{m}$, $1\,\mathrm{km}$), temporal resolution towers
  (daily, monthly, annual), and semantic resolution towers
  (species, genus, family), and to systems where all three
  hierarchies are active simultaneously.
\end{remark}


% ============================================================
\begin{thebibliography}{99}

\bibitem{Flyxion2026a}
Flyxion.
\newblock Observability is restorability: Restorable distinctions and the
  geometry of perceptual access.
\newblock Preprint, 2026.

\bibitem{Flyxion2026b}
Flyxion.
\newblock Erasure is deformation: Repair-repertoire reduction and the
  geometry of memory loss.
\newblock Preprint, 2026.

\bibitem{Weikmann2026}
Weikmann, Giulio, Gianmarco Perantoni, and Lorenzo Bruzzone.
\newblock A semantics-aware hierarchical consensus approach to satellite image
  time series classification.
\newblock \emph{arXiv}:2510.04916, under review for \emph{IEEE Transactions
  on Geoscience and Remote Sensing}, 2026.

\bibitem{Kuhn1962}
Kuhn, Thomas~S.
\newblock \emph{The Structure of Scientific Revolutions}.
\newblock University of Chicago Press, Chicago, 1962.

\bibitem{Kleene1952}
Kleene, Stephen~C.
\newblock \emph{Introduction to Metamathematics}.
\newblock North-Holland, Amsterdam, 1952.

\bibitem{Tarski1944}
Tarski, Alfred.
\newblock The semantic conception of truth and the foundations of semantics.
\newblock \emph{Philosophy and Phenomenological Research}, 4(3):341--376, 1944.

\bibitem{Frege1892}
Frege, Gottlob.
\newblock \"{U}ber Sinn und Bedeutung.
\newblock \emph{Zeitschrift f\"{u}r Philosophie und philosophische Kritik},
  100:25--50, 1892.

\bibitem{Quine1951}
Quine, Willard~Van~Orman.
\newblock Two dogmas of empiricism.
\newblock \emph{The Philosophical Review}, 60(1):20--43, 1951.

\bibitem{Wittgenstein1921}
Wittgenstein, Ludwig.
\newblock \emph{Tractatus Logico-Philosophicus}.
\newblock Kegan Paul, Trench, Trubner \& Co., London, 1922.

\bibitem{Munkres2000}
Munkres, James~R.
\newblock \emph{Topology}.
\newblock Prentice Hall, Upper Saddle River, NJ, 2nd edition, 2000.

\end{thebibliography}

\end{document}
